\documentclass[11pt]{article}
\usepackage[margin=1in]{geometry}
\usepackage[T1]{fontenc}
\usepackage[utf8]{inputenc}
\usepackage{lmodern}
\usepackage{microtype}
\usepackage{amsmath,amssymb,amsthm,mathtools}
\usepackage{enumitem}
\usepackage{booktabs}
\usepackage{array}
\usepackage{longtable}
\usepackage{hyperref}
\usepackage{xcolor}
\hypersetup{colorlinks=true,linkcolor=blue,citecolor=blue,urlcolor=blue}
\setlist[itemize]{leftmargin=1.5em}
\setlist[enumerate]{leftmargin=1.75em}

\title{A Conceptual Primer on World-Model Markov Reasoning,\\Hidden-State Inference, and POMDP Belief Semantics\\\large Dated synthesis: 2026-04-17}
\author{}
\date{}

\newtheorem{definition}{Definition}
\newtheorem{principle}{Principle}
\newtheorem{claim}{Claim}
\newtheorem{remark}{Remark}
\newcommand{\wm}{world-model}
\newcommand{\wmpln}{WM-PLN}
\newcommand{\pln}{PLN}
\newcommand{\R}{\mathbb{R}}
\newcommand{\ENN}{\mathbb{R}_{\ge 0}^{\infty}}
\newcommand{\argmax}{\mathop{\mathrm{argmax}}}
\newcommand{\Law}{\mathrm{Law}}
\newcommand{\Bel}{\mathrm{Bel}}
\newcommand{\Obs}{\mathrm{Obs}}

\begin{document}
\maketitle

\begin{abstract}
This note is a conceptual primer for a research line connecting world-model
calculus, Markov semantics, higher-order Markov chains, hidden Markov models,
credal uncertainty, and POMDP-style reasoning.  Its purpose is not to archive
implementation history but to preserve the strongest conceptual picture now
available so that future technical work and future conversations can resume
from a stable basis.

The central thesis is that probabilistic symbolic reasoning becomes reliable
only when it sharply separates five layers that are often conflated:
\begin{enumerate}
  \item the explicit semantic model that determines probabilities,
  \item the evidence or posterior state carried by a reasoning system,
  \item the query surface exposed to users,
  \item the operational control layer that decides what computation to perform,
  \item and the higher-order uncertainty induced by hidden structure,
        non-identifiability, or model ambiguity.
\end{enumerate}

On this view, additive evidence revision remains a mathematically privileged
fragment, but it is not the universal semantics of hidden-state reasoning.  For
ordinary finite-state Markov chains, predictive path semantics can be exposed
honestly through world-model queries.  For higher-order Markov chains, the
right reduction is to first-order chains on context states.  For hidden Markov
models, complete-data evidence remains additive, but observed-only belief
tracking is fundamentally non-additive; the correct response is filtering,
smoothing, and, in the ambiguous case, credal or interval-valued belief
representations.  Action-conditioned hidden Markov models then provide a clean
structured entry point for POMDP-capable world-model reasoning.
\end{abstract}

\tableofcontents

\section{Purpose and scope}

This document consolidates a cluster of ideas that belong together:
\begin{itemize}
  \item world-model semantics as the truthmaking layer for probabilistic logic;
  \item honest Markov path semantics in place of ad hoc multi-step chaining;
  \item higher-order Markov semantics via context-state reduction;
  \item hidden-state inference via forward, backward, filtering, and smoothing;
  \item observed-only non-identifiability and its implications for additive
        evidence calculi;
  \item credal world-model views as an honest response to ambiguity;
  \item and controlled hidden Markov models as a bridge into POMDP-style
        reasoning.
\end{itemize}

The note is strictly conceptual.  It is written to be useful both as an
academic synthesis and as a primer for future technical work.  It therefore
emphasizes conceptual boundaries, canonical constructions, and theorem-shaped
open problems rather than implementation detail.

\section{The guiding separation: semantics, evidence, queries, and control}

\subsection{The main distinction}

The first indispensable distinction is:
\begin{quote}
The semantic model determines what is true; the carried evidence or posterior
state records what a reasoner knows; the query surface exposes selected views of
that state; and the control layer decides what computation to perform next.
\end{quote}

Many failures of probabilistic symbolic reasoning arise from collapsing these
roles into one object.

\begin{itemize}
  \item A probability in a semantic model is not the same thing as a confidence
        attached to a derived symbolic judgment.
  \item A local truth value is not automatically a reusable semantic fact.
  \item A search-control score is not the same thing as posterior support.
  \item A posterior belief state is not the same thing as an additive multiset
        of observations unless the underlying model class genuinely supports
        such a reduction.
\end{itemize}

\subsection{The additive fragment remains special}

The additive \wm{} fragment should not be discarded.  It remains canonical for
families where independent observations aggregate through sufficient
statistics.  In that setting, evidence accumulation is literally additive.  The
point is not that additivity is wrong; the point is that it is correct only in
the fragment where the semantics really supports it.

\begin{principle}[Additive privilege]
When a model class admits a sufficient-statistic posterior that aggregates by
addition, the additive world-model semantics is the canonical surface for that
class.
\end{principle}

Positive example:
\begin{itemize}
  \item Beta--Bernoulli and Dirichlet--multinomial families aggregate evidence
        by counts, so additive evidence states are exactly right.
\end{itemize}

Negative example:
\begin{itemize}
  \item Observed-only hidden-state inference generally does not reduce to an
        additive multiset statistic that recovers the correct posterior over
        latent states.
\end{itemize}

\section{From local chaining to explicit predictive semantics}

\subsection{Why naive probabilistic chaining fails}

The tempting but unreliable idea is that one can propagate local truth values
through arbitrary chains as if local support values composed like logical
probabilities.  This fails for familiar reasons:
\begin{itemize}
  \item hidden overlap and dependence are forgotten,
  \item derived links are reused as if they were independent raw premises,
  \item and multi-step composition is forced into an additive or deductive form
        that the semantics does not justify.
\end{itemize}

\subsection{What survives}

What survives is the narrower but much cleaner claim:
\begin{quote}
Multi-step probabilistic reasoning should be exposed through direct semantic
queries against an explicit model, not through generic reuse of locally derived
support values.
\end{quote}

For finite-state Markov chains, this means that multi-step queries should be
read as \emph{predictive path probabilities}.  The right semantics is
sequential posterior or predictive composition, not generic additive deduction.

\section{Finite-state Markov world-model semantics}

\subsection{Single-step queries}

The simplest Markov query asks for the probability of a one-step transition
\[
P(X_{t+1}=x' \mid H_t),
\]
where $H_t$ is the relevant history summary.  In a conjugate finite-state
setting, the sufficient statistic can be read off from row counts and boundary
data.

\subsection{Path queries}

The conceptual improvement over naive chaining is to introduce direct path
queries
\[
P(X_{t+1}=x_1,\ldots,X_{t+n}=x_n \mid H_t).
\]
These are not mere chains of local truth values; they are model-defined path
probabilities.

\begin{principle}[Predictive over deductive]
For finite-state Markov reasoning, multi-step query strength should be defined
by predictive path semantics, not by additive reuse of locally derived
transition judgments.
\end{principle}

\subsection{Threshold semantics}

Once direct predictive query strength exists, threshold theorems become
meaningful: one may ask whether a path or transition has support above a given
cutoff, but this thresholding happens \emph{after} semantic evaluation, not in
place of it.

\section{Higher-order Markov chains through context-state reduction}

\subsection{The basic idea}

An order-$m$ Markov chain on symbols can be reinterpreted as a first-order
Markov chain on length-$m$ contexts.  This is not merely a coding trick.  It
is the mathematically right reduction for bringing higher-order chains into the
same finite-state kernel and world-model infrastructure.

\begin{definition}[Context-state reduction]
An order-$m$ chain on symbols $x_t$ is represented as a first-order chain on
states
\[
c_t = (x_t,\ldots,x_{t+m-1}),
\]
with transitions given by shift-consistent context extension.
\end{definition}

\subsection{Why this matters}

This reduction gives three conceptual payoffs.
\begin{enumerate}
  \item It turns higher-order sequence laws into ordinary first-order path laws
        on context states.
  \item It identifies the right sufficient statistic: $(m+1)$-gram counts.
  \item It makes higher-order reasoning compatible with the same predictive and
        evidence-carrying surfaces used for first-order finite-state chains.
\end{enumerate}

\subsection{Raw symbol semantics versus context semantics}

There are two distinct but equivalent views:
\begin{itemize}
  \item context-trajectory semantics, where the underlying process lives on
        contexts;
  \item raw-symbol semantics, where one projects context trajectories back down
        to symbol sequences.
\end{itemize}

The raw-symbol cylinder law is conceptually important because it shows that the
context-state reduction is not just an auxiliary proof device; it is genuinely
equivalent to the higher-order symbolic process.

\section{Hidden Markov models: from complete data to observed-only inference}

\subsection{The two views of an HMM}

A finite HMM has:
\begin{itemize}
  \item a latent Markov chain,
  \item an emission kernel,
  \item and an induced observed process obtained by summing or integrating over
        latent paths.
\end{itemize}

From the complete-data point of view, the model remains additive in a useful
way: one can accumulate latent transition counts and emission counts.  From the
observed-only point of view, however, the correct object is not an additive
count summary but a posterior distribution over latent states.

\subsection{Forward and backward semantics}

The forward message computes the unnormalized mass of ending in latent state
$x$ after an observed prefix.  The backward message computes the continuation
likelihood of the future observations given latent state $x$.  Their product
recovers the standard smoothing factorization.

Conceptually, four objects matter:
\begin{itemize}
  \item observed-word probability,
  \item filtering posterior,
  \item smoothing posterior,
  \item and two-slice latent posterior.
\end{itemize}

\subsection{Filtering}

Filtering is the posterior belief state
\[
\Bel_t(x) = P(X_t=x \mid Y_{\le t}, A_{\le t}),
\]
or its action-free specialization when no actions are present.  This is the
load-bearing object for observed-only reasoning.

\subsection{Smoothing}

Smoothing is the two-sided posterior
\[
P(X_t=x \mid Y_{\le T}, A_{\le T-1}),
\]
or its finite-prefix analogue.  The conceptual lesson is that smoothing is not
an optional embellishment.  It makes explicit that observed-only reasoning is
posterior reasoning over hidden trajectories, not a direct additive tally.

\section{The non-identifiability boundary}

\subsection{What the obstruction says}

The central negative insight is:
\begin{quote}
Observed-only hidden-state belief cannot in general be recovered from an
additive multiset-based evidence surface over observations alone.
\end{quote}

The cleanest witness is a label-swap construction: two hidden models can induce
the same observed distribution while assigning incompatible latent posteriors.
This is not a pathology.  It is the structural reason observed-only HMM belief
tracking cannot simply inherit the complete-data additive world-model
infrastructure.

\subsection{Why the result matters}

This obstruction has three consequences.
\begin{enumerate}
  \item It marks the boundary of additive \wmpln{} reasoning for hidden-state
        models.
  \item It justifies filtering and smoothing as first-class semantic objects.
  \item It motivates credal or interval-valued belief representations when the
        observations underdetermine latent structure.
\end{enumerate}

\section{Credal world-model semantics for observed-only ambiguity}

\subsection{The conceptual move}

When the observations do not identify a unique hidden model, the right response
is not to pretend that a single posterior strength exists.  Instead one should
carry a family of compatible hidden models and expose lower and upper support
for a latent query.

\begin{definition}[Credal filtering truth value]
Given a finite family of compatible hidden models $\Theta$, an observed trace
$z$, and a latent query $q$, define
\[
\underline{s}(q \mid z) = \inf_{\theta \in \Theta} P_\theta(q \mid z), \qquad
\overline{s}(q \mid z) = \sup_{\theta \in \Theta} P_\theta(q \mid z).
\]
\end{definition}

This interval is the honest observed-only world-model answer when latent
structure is ambiguous.

\subsection{Why this is conceptually natural}

The credal move is not a retreat from semantics.  It is the semantics of
ambiguity.  In this sense, the obstruction theorem is not merely a negative
result.  It says what the next representation should be.

Positive example:
\begin{itemize}
  \item a label-swapped HMM family can induce the full interval $[0,1]$ for a
        latent-state query after a singleton observation.
\end{itemize}

Negative example:
\begin{itemize}
  \item the interval need not collapse even for very short traces, so
        observed-only ambiguity is not a rare asymptotic phenomenon.
\end{itemize}

\section{Controlled HMMs as a structured bridge to POMDP reasoning}

\subsection{Why action-conditioning matters}

A POMDP is not merely an HMM with observations.  It is an action-conditioned
partially observable process.  The clean structured bridge is therefore a
controlled finite HMM:
\begin{itemize}
  \item latent prior,
  \item action-conditioned latent transition kernel,
  \item emission kernel,
  \item filtering and smoothing on completed action-observation cycles.
\end{itemize}

\subsection{What this unlocks}

This does \emph{not} create the general agent framework from scratch.  Rather,
it provides a structured finite-state instance of an already general
history-conditional environment formalism.  The main gain is conceptual:
\begin{quote}
The belief state of a finite POMDP can now be expressed as a controlled
filtering posterior in the world-model reasoning layer.
\end{quote}

\subsection{The controlled Solomonoff leg}

The controlled story now has a conservative universal-prediction leg as well.
A controlled finite HMM induces a controlled prefix law on completed
action-observation traces, and along any fixed action stream one can condition
that law down to an ordinary observation-prefix measure.  The ordinary
finite-alphabet Solomonoff semimeasure $M_2$ can then be lifted back into the
controlled setting by ignoring actions and acting as a comparator on the
observation side.  This yields a real theorem:
\begin{quote}
For any fixed action stream, if the induced observation law is lower
semicomputable, then the controlled Solomonoff comparator satisfies the
standard finite-horizon log-loss regret bound against that controlled hidden
Markov law.
\end{quote}

Positive example:
\begin{itemize}
  \item the controlled Markov / Solomonoff / \wm{} triangle is now closed at
        the level of fixed-action-stream regret bounds.
\end{itemize}

Negative example:
\begin{itemize}
  \item the current comparator is deliberately conservative: it is not yet a
        genuinely action-aware universal mixture.
\end{itemize}

\subsection{A genuine computability boundary}

The lower-semicomputability side condition in the controlled Solomonoff theorem
is not a technical embarrassment.  It reflects a real domain boundary.  The
raw controlled finite HMM parameter record allows arbitrary probability
measures, so lower semicomputability does not follow automatically from the
bare semantic object.  The right response is to package a computable or
lower-semicomputable subclass explicitly, not to erase the assumption by fiat.

\subsection{The correct interpretation}

In this context, ``unlocking POMDP reasoning'' means:
\begin{itemize}
  \item one has a structured controlled latent-state model,
  \item one has a bridge from that model to a generic environment interface,
  \item and one can state planning semantics in terms of belief-state update.
\end{itemize}

What it does \emph{not} mean is that additive evidence alone now suffices for
hidden-state planning.  The obstruction theorem still applies.

\section{World-model reasoning for POMDPs}

\subsection{What proper WM reasoning should mean here}

Proper \wm{} reasoning for POMDP-like systems should mean:
\begin{enumerate}
  \item a semantic model with explicit latent and observed dynamics,
  \item belief-state update equations,
  \item query surfaces for filtering, smoothing, and prediction,
  \item planning or control semantics that read those belief states correctly,
  \item and an honest ambiguity layer when observations fail to identify a
        unique latent state model.
\end{enumerate}

\subsection{What it should not mean}

It should not mean:
\begin{itemize}
  \item reusing additive observation multisets as if they were sufficient for
        latent-state posteriors,
  \item calling any local truth-value propagation scheme ``belief update'',
  \item or hiding hidden-state ambiguity under a single scalar confidence.
\end{itemize}

\section{The de Finetti and mixture story}

\subsection{What is trivial and what is not}

Mixture theorems come in two sharply different forms.

\paragraph{Trivial direction.}
For each fixed parameter $\theta$, the law induced by $\theta$ can always be
represented as a Dirac mixture concentrated at $\theta$.

\paragraph{Nontrivial direction.}
One wants a real mixture theorem: finite discrete mixtures, or ultimately
arbitrary Borel mixtures, together with an honest image characterization of the
laws thus obtained.

\subsection{The naming lesson}

The conceptual lesson is methodological:
\begin{quote}
Dirac-witness existence theorems are legitimate, but they should be named as
Dirac-witness theorems rather than as if they already established the full de
Finetti image theorem.
\end{quote}

This matters because the difference is mathematically substantive.  A statement
may be correct as written while still overselling what kind of factorization has
actually been proved.

\subsection{Higher-order and HMM mixture surfaces}

For higher-order Markov models, the first honest nontrivial step is a finite
discrete mixture theorem on context-state parameters.  For hidden Markov
models, the same applies to latent Markov parameters with fixed emission
kernel.  The genuinely hard open layer is the arbitrary-Borel image theorem in
both settings.

\section{The ProbMeTTa lesson: semantic mirror versus source semantics}

\subsection{The initial gap}

Another general lesson emerging from this line of work is the distinction
between:
\begin{itemize}
  \item a semantic theorem downstream of a compilation or translation relation,
  \item and a formal semantics of the actual source-level runtime language.
\end{itemize}

It is easy to prove a theorem of the form:
\begin{quote}
If the source system compiles to a certain semantic object, then the downstream
theorem holds.
\end{quote}
This is useful, but it is not the same as formalizing the actual source
operators and proving the compiler relation from the runtime semantics.

\subsection{The right target}

The right target is not byte-for-byte replay of an implementation.  The right
target is a mathematically honest source semantics for the runtime core,
together with adequacy theorems connecting it to the already-proved semantic
backend.

This yields a more general methodological rule:
\begin{principle}[Semantic fidelity over operational mimicry]
When formalizing a symbolic runtime, the goal should be semantic adequacy of the
source operators, not a brittle proof of implementation-level reduction order.
\end{principle}

\section{Stable design rules for future work}

The following rules are worth treating as stable design commitments.

\begin{enumerate}
  \item \textbf{Name theorems honestly.}  Distinguish Dirac witnesses from real
        mixture theorems.
  \item \textbf{Do not publish consumer surfaces before the semantics are
        stable.}  Query syntax should follow semantic stabilization, not
        precede it.
  \item \textbf{Prefer direct semantic queries over reuse of derived links.}
  \item \textbf{Treat additive evidence as a privileged fragment, not a
        universal law.}
  \item \textbf{Use credal or interval-valued views when hidden structure is
        not identifiable from observations.}
  \item \textbf{Keep complete-data and observed-only semantics separate.}
  \item \textbf{Reduce higher-order chains through context states unless a
        better semantic representation is proved.}
  \item \textbf{Bridge into planning through controlled filtering, not through
        ad hoc symbolic chaining.}
\end{enumerate}

\section{The next theorem program}

The most coherent next theorem program is:
\begin{enumerate}
  \item a computable or lower-semicomputable controlled-HMM parameter class
        that discharges the fixed-action-stream Solomonoff hypothesis by
        construction;
  \item stronger credal planning bounds beyond the current finite-horizon
        envelope theorems;
  \item ambiguity-sensitive planning examples where observation-consistent model
        families widen the decision interval, not only the posterior interval;
  \item arbitrary-Borel mixture theorems for higher-order Markov and hidden
        Markov settings;
  \item a genuinely action-aware universal mixture, rather than the current
        conservative action-ignoring comparator;
  \item and only after those, an HMM/POMDP query surface for symbolic consumer
        interfaces.
\end{enumerate}

\section{A compact primer for future technical conversations}

Future work can safely assume the following conceptual baseline.

\begin{enumerate}
  \item \textbf{Finite-state Markov chains:} direct path semantics is the right
        multi-step query surface.
  \item \textbf{Higher-order Markov chains:} the right reduction is to
        first-order chains on context states; the sufficient statistics are
        $(m+1)$-gram counts.
  \item \textbf{Hidden Markov models:} complete-data evidence is additive, but
        observed-only latent belief is posterior, not additive tally.
  \item \textbf{Observed-only ambiguity:} when multiple hidden models remain
        compatible with the same observations, lower and upper latent support
        should be carried explicitly.
  \item \textbf{POMDP bridge:} the structured finite-state route goes through
        controlled hidden Markov models and belief-state recursion.
  \item \textbf{Controlled Solomonoff leg:} along fixed action streams there is
        now a concrete conservative universal comparator with a genuine
        finite-horizon regret theorem.
  \item \textbf{Formalization strategy:} prefer semantic adequacy and honest
        boundaries over overclaimed universality or implementation mimicry.
  \item \textbf{Open boundary:} the arbitrary-Borel mixture and image
        theorems remain the deepest unresolved measure-theoretic layer.
  \item \textbf{Practical implication:} proper symbolic probabilistic reasoning
        requires different surfaces for additive evidence, predictive path
        queries, posterior latent beliefs, and ambiguity-aware credal views.
\end{enumerate}

\section{Conclusion}

The broad picture is now reasonably clear.

The right semantics for probabilistic symbolic reasoning is not generic
truth-value propagation.  It is world-model semantics with query-specific
evidence extraction.  Additive evidence remains canonical where sufficient
statistics justify it.  Predictive path semantics is the right surface for
Markov composition.  Higher-order chains should be reduced through context
states.  Hidden-state models require filtering and smoothing rather than naive
observation counting.  Observed-only ambiguity calls for credal world-model
views.  Action-conditioned hidden Markov models provide the cleanest bridge to
POMDP-capable world-model reasoning.

This yields a conceptually disciplined research program: one preserves the
fragments that are exactly additive, one introduces posterior objects where
additivity genuinely fails, one names theorems honestly, and one delays
consumer surfaces until the semantic core has stabilized.  That is a much
cleaner foundation for future reasoning systems than either naive local
probabilistic chaining or overgeneralized confidence calculus.

\end{document}
